\documentclass[conference]{IEEEtran}
\IEEEoverridecommandlockouts
% The preceding line is only needed to identify funding in the first footnote. If that is unneeded, please comment it out.
\usepackage{cite}
\usepackage{amsmath,amssymb,amsfonts}
\usepackage{algorithmic}
\usepackage{graphicx}
\usepackage{textcomp}
\usepackage{xcolor}
\usepackage{booktabs}
\usepackage{url}
\usepackage{multirow}
\usepackage{array}

\def\BibTeX{{\rm B\kern-.05em{\sc i\kern-.025em b}\kern-.08em
    T\kern-.1667em\lower.7ex\hbox{E}\kern-.125emX}}

\begin{document}

\title{Architecture, Elasticity, and Automated Delivery of an Event-Driven Video Transcoding Platform on Serverless Container Infrastructure}

\author{
\IEEEauthorblockN{Shlok Doshi}
\IEEEauthorblockA{\textit{Roll No: B039} \\
\textit{Dept. of Computer Engineering} \\
Mumbai, India}
\and
\IEEEauthorblockN{Kartika Dhonde}
\IEEEauthorblockA{\textit{Roll No: B037} \\
\textit{Dept. of Computer Engineering} \\
Mumbai, India}
\and
\IEEEauthorblockN{Samar Kamat}
\IEEEauthorblockA{\textit{Roll No: B059} \\
\textit{Dept. of Computer Engineering} \\
Mumbai, India}
\and
\IEEEauthorblockN{Kshitij Jhalak}
\IEEEauthorblockA{\textit{Roll No: B055} \\
\textit{Dept. of Computer Engineering} \\
Mumbai, India}
}

\maketitle

\begin{abstract}
Modern multimedia applications demand highly scalable, resilient, and cost-efficient processing pipelines to handle unpredictable bursts of user-generated content. However, designing such platforms introduces complex architectural trade-offs among compute elasticity, operational simplicity, release safety, and cloud expenditure. This paper presents the architecture, end-to-end implementation, and comprehensive empirical evaluation of a cloud-native, event-driven video upload, transcoding, and streaming platform deployed on serverless container infrastructure using AWS Elastic Container Service (ECS) on AWS Fargate and Fargate Spot. 

The system decouples ingestion, batch processing, and global content delivery: (1) high-throughput ingestion is achieved via client-direct pre-signed Amazon S3 uploads managed by a lightweight Python/FastAPI service, eliminating intermediate container memory buffering; (2) asynchronous job scheduling is orchestrated via Amazon S3 event notifications, Amazon SQS message queues, and Amazon DynamoDB for state tracking; (3) compute-intensive multi-rendition transcoding (720p, 480p H.264/AAC) and thumbnail extraction are executed by FFmpeg worker containers running on cost-optimized AWS Fargate Spot instances; (4) worker provisioning leverages queue-depth-driven step auto-scaling, scaling dynamically from zero to three instances to achieve near-zero idle compute cost; (5) delivery safety is enforced via GitHub Actions continuous integration and continuous deployment (CI/CD) pipelines utilizing OpenID Connect (OIDC), Trivy container vulnerability scanning, Amazon ECS deployment circuit breakers with automated rollback, and SQS Dead-Letter Queue (DLQ) redrive mechanisms. 

We evaluate the system through an extensive benchmark suite simulating production workloads across standardized 1080p video clips. Experimental results reveal a 69.8\% compute cost reduction via Fargate Spot compared to on-demand compute, an average turnaround latency of 38.1 seconds for 30-second source clips, automated circuit-breaker rollback recovery in 84.6 seconds, and an 82.4\% operational cost reduction compared to commercial managed media transcoding services.
\end{abstract}

\begin{IEEEkeywords}
Cloud Computing, Microservices, Video Transcoding, Serverless Containers, AWS Fargate Spot, Event-Driven Architecture, Auto-Scaling, Continuous Integration/Continuous Deployment (CI/CD), Deployment Circuit Breaker, Fault Tolerance.
\end{IEEEkeywords}

\section{Introduction}
Online video consumption accounts for over 82\% of global internet traffic, driven by social platforms, video sharing portals, e-learning environments, and teleconferencing \cite{cisco}. Processing user-uploaded media is computationally demanding: source videos feature heterogeneous codecs, bitrates, resolutions, and container formats. To guarantee uniform playback across heterogeneous client devices and fluctuating network bandwidths, platforms must transcode raw video streams into multiple standardized renditions and compress them for efficient content delivery network (CDN) distribution \cite{timmerer}.

Video transcoding is an embarrassingly parallel, compute-bound, and highly bursty workload. Traditional monolithic architectures—where transcoding runs synchronously on statically provisioned virtual machines (VMs)—exhibit major operational inefficiencies:
\begin{itemize}
    \item \textbf{Resource Inefficiency:} Static VM provisioning incurs substantial financial overhead during traffic troughs and severe queue buildup or dropped connections during load spikes \cite{armbrust}.
    \item \textbf{Ingress Bottlenecks:} Routing large multimedia payloads through intermediate API gateways or web servers saturates container memory buffers, exhausts network socket pools, and starves concurrent control-plane requests \cite{nygard}.
    \item \textbf{Deployment Fragility:} Monolithic applications carry high blast radiuses; an unhandled exception or native dependency collision during transcoding can crash user-facing metadata and playback endpoints \cite{richardson}.
\end{itemize}

Cloud-native paradigms advocate for microservice decomposition, containerization, and event-driven asynchronous messaging \cite{dragoni, jamshidi}. Serverless container engines, such as AWS Fargate, allow developers to run containers without managing the underlying virtual machine clusters \cite{burns}. Furthermore, surplus cloud capacity offerings (e.g., AWS Fargate Spot) offer discounts up to 70\% for fault-tolerant batch workloads \cite{fargate_spot}.

Despite these theoretical advantages, small-to-medium engineering teams face significant architectural challenges in synthesizing these primitives into a unified, resilient, and economically sustainable platform. Key engineering hurdles include: (i) mitigating scale-from-zero cold start delays for bursty queues; (ii) guaranteeing automated release rollbacks without operator intervention; and (iii) eliminating high-cost cloud components (e.g., NAT Gateways) without compromising container security.

\subsection{Contributions}
This paper presents the design, implementation, and empirical evaluation of an open-source, cloud-native video upload, transcoding, and streaming platform. The core contributions of this work are:
\begin{enumerate}
    \item \textbf{Decoupled Event-Driven Reference Architecture:} We formulate a microservices architecture combining FastAPI, Amazon S3, SQS, DynamoDB, ECS Fargate Spot, and CloudFront with Origin Access Control (OAC), eliminating NAT Gateways through public-subnet security isolation.
    \item \textbf{Direct-to-Storage Pre-signed Ingress Pattern:} We analyze the memory and throughput impact of client-direct S3 pre-signed uploads versus intermediate server proxying, demonstrating a 99.4\% reduction in API memory overhead.
    \item \textbf{Queue-Depth Step Auto-Scaling with Scale-to-Zero:} We evaluate a step scaling policy driven by Amazon CloudWatch queue-backlog metrics that provisions Fargate Spot worker tasks while gracefully handling POSIX \texttt{SIGTERM} signals during Spot evictions.
    \item \textbf{Dual-Layer Rollback \& Release Safety Framework:} We implement and validate an automated CI/CD pipeline incorporating OIDC passwordless authentication, Aquasecurity Trivy CVE scanning, ECS deployment circuit breakers, and SQS dead-letter queue (DLQ) redrive routines.
    \item \textbf{Empirical Benchmarking \& Economic Modeling:} We present an extensive experimental evaluation covering transcoding turnaround latency, CPU/memory profiles, auto-scaling cold starts, CI/CD deployment lead times, and multi-model cloud cost comparisons.
\end{enumerate}

\section{Background and Related Work}

\subsection{Containerization and Lightweight Virtualization}
Operating-system-level virtualization through Linux containers encapsulates application binaries and dependencies into immutable images, guaranteeing runtime parity across development and cloud environments \cite{merkel}. Unlike hypervisor-based virtual machines that duplicate entire guest operating systems, containers share the host Linux kernel, resulting in minimal memory overhead and rapid process instantiation \cite{felter}.

Managed container orchestrators have evolved from complex self-hosted Kubernetes clusters to serverless container runtimes such as AWS ECS with AWS Fargate \cite{burns}. AWS Fargate isolates container tasks using Firecracker micro-virtual machines—lightweight software execution environments built on Linux Kernel-based Virtual Machine (KVM) technology that provide hypervisor-grade isolation with sub-second boot times and minimal memory footprints \cite{agache}. Fargate Spot allows tasks to run on unused compute capacity at substantial price discounts, subject to occasional two-minute interruption notices.

\subsection{Microservices and Asynchronous Messaging}
The microservices architectural pattern organizes systems into loosely coupled, independently scalable services centered around discrete business capabilities \cite{fowler, newman}. To avoid synchronous cascading failures across distributed microservices, asynchronous message queues are widely employed \cite{taibi}.

Hohpe and Woolf formalized core enterprise integration patterns, including \textit{Competing Consumers}, \textit{Dead-Letter Channels}, and \textit{Idempotent Receivers} \cite{hohpe}. In a competing consumers topology, worker processes poll a durable message queue (such as Amazon SQS), absorbing traffic surges and leveling compute load. If a worker encounters an unhandled exception or terminates prematurely, the message visibility timeout expires, returning the message to the queue for reprocessing. If retries exceed a configured threshold ($M_{\text{recv}} = 3$), the message is diverted to a Dead-Letter Queue (DLQ), preserving system stability and enabling post-mortem diagnosis.

\subsection{Cloud Video Transcoding Architectures}
Video transcoding transforms source bitstreams into target formats through decoding, scaling, and re-encoding \cite{wiegand}. Research into cloud-based video transcoding generally falls into two paradigms:
\begin{itemize}
    \item \textbf{Fine-Grained Serverless Transcoding:} Frameworks like ExCamera \cite{fouladi} and Sprocket \cite{ao} leverage thousands of ephemeral AWS Lambda functions to transcode small Groups of Pictures (GOPs) in parallel. While achieving near-real-time latency, fine-grained serverless encoding incurs significant coordination complexity, complex audio/video re-muxing overhead, and high per-request S3 API costs.
    \item \textbf{Containerized Batch Transcoding:} Platforms utilizing containerized worker daemons running FFmpeg \cite{bellard} process discrete video files sequentially. Jokhio et al. \cite{jokhio} demonstrated that containerized transcoders offer predictable resource utilization and superior cost efficiency for asynchronous media processing. Our platform adopts this model while extending it with scale-to-zero Fargate Spot elasticity.
\end{itemize}

\subsection{Continuous Delivery and Release Engineering}
Continuous Delivery (CD) automates the validation and promotion of software releases \cite{humble}. Forsgren et al. established that delivery velocity and release stability (measured via deployment frequency, lead time for changes, change failure rate, and mean time to restore service) directly predict operational performance \cite{forsgren}. Modern deployment orchestrators implement \textit{Deployment Circuit Breakers} \cite{ecs_breaker}, which actively monitor downstream health check telemetry during rolling updates. When newly launched container revisions fail consecutive health checks, the circuit breaker automatically halts the rollout and rolls back service configurations to the prior healthy revision ($N-1$).

\section{System Architecture and Implementation}

The platform architecture is designed around four foundational principles: **stateless compute execution**, **storage-level ingress decoupling**, **event-driven choreography**, and **infrastructure immutability**. 

\begin{figure*}[t]
\centering
\fbox{
\begin{minipage}{0.96\textwidth}
\small
\textbf{[Client Browser]} $\xrightarrow{\text{1. POST /api/videos}}$ \textbf{[ALB]} $\rightarrow$ \textbf{[FastAPI on ECS Fargate]} $\rightarrow$ \textbf{[DynamoDB (status="uploaded")]} \\
\hspace*{2.4cm} $\xleftarrow{\text{2. Pre-signed S3 PUT URL}}$ \\
\textbf{[Client Browser]} $\xrightarrow{\text{3. Direct PUT raw/\{id\}.mp4}}$ \textbf{[S3 Uploads Bucket]} \\
\hspace*{4.8cm} $\xrightarrow{\text{4. s3:ObjectCreated}}$ \textbf{[SQS video-jobs]} $\rightarrow$ \textit{(CloudWatch Alarm scales Fargate Spot Workers 0..3)} \\
\hspace*{8.2cm} $\xrightarrow{\text{5. SQS Poll}}$ \textbf{[ECS Fargate Spot Worker]} \\
\hspace*{10.8cm} $\rightarrow$ FFmpeg 720p / 480p / Thumb \\
\hspace*{10.8cm} $\rightarrow$ Upload to \textbf{[S3 Processed Bucket]} \\
\hspace*{10.8cm} $\rightarrow$ Update \textbf{[DynamoDB (status="ready")]} \\
\textbf{[Client Browser]} $\xleftarrow{\text{6. HTTPS Playback /processed/*}}$ \textbf{[Amazon CloudFront (Origin Access Control)]} $\leftarrow$ \textbf{[S3 Processed Bucket]}
\end{minipage}
}
\caption{End-to-end event-driven video ingestion, transcoding, and content delivery pipeline.}
\label{fig:architecture}
\end{figure*}

\subsection{Direct-to-Storage Pre-signed Ingress Pattern}
Routing video uploads through application containers saturates web worker threads and exhausts network buffers. To eliminate intermediate ingress bottlenecks, the platform implements a direct-to-storage pre-signed ingress pattern:
\begin{enumerate}
    \item The client initiates an upload request by sending metadata (\texttt{\{"title": "Demo"\}}) to the FastAPI service (\texttt{POST /api/videos}).
    \item The API generates a unique Version 4 UUID (\texttt{video\_id}) and writes an initial item to Amazon DynamoDB with \texttt{status = "uploaded"} and \texttt{source\_key = "raw/\{video\_id\}.mp4"}.
    \item The API invokes the AWS Boto3 SDK to compute an HMAC-SHA256 pre-signed Amazon S3 PUT URL valid for 900 seconds:
    \begin{equation}
        \text{URL} = \text{HMAC-SHA256}(\mathcal{K}_{\text{secret}}, \mathcal{P}_{\text{request}}, \mathcal{T}_{\text{stamp}})
    \end{equation}
    \item The client receives the signature and transmits the raw MP4 file directly to the S3 \texttt{uploads} bucket via HTTP PUT. The API container never buffers or routes media bytes.
\end{enumerate}

\subsection{Asynchronous Event Dispatch and Distributed State Machine}
When the upload to S3 completes, Amazon S3 generates an \texttt{s3:ObjectCreated:Put} event notification, which is immediately dispatched into an Amazon SQS standard message queue (\texttt{video-jobs}). 

A distributed state machine manages each video lifecycle in Amazon DynamoDB:
\begin{equation}
    \mathcal{S} \in \{\text{uploaded}, \text{processing}, \text{ready}, \text{failed}\}
\end{equation}
Table \ref{tab:dynamodb} details the data contract enforced across the microservices.

\begin{table}[htbp]
\caption{DynamoDB Schema and Data Contract (\texttt{video-platform-videos})}
\label{tab:dynamodb}
\begin{center}
\begin{tabular}{llll}
\toprule
\textbf{Attribute} & \textbf{Type} & \textbf{Mutability} & \textbf{Description} \\
\midrule
\texttt{video\_id} & String (PK) & Immutable & UUIDv4 primary identifier \\
\texttt{title} & String & Immutable & User-provided title \\
\texttt{status} & String & Mutable & State: uploaded, processing, ready, failed \\
\texttt{source\_key} & String & Immutable & S3 path of raw input video \\
\texttt{outputs} & Map & Mutable & CloudFront paths for 720p/480p MP4s \\
\texttt{thumbnail\_key} & String & Mutable & S3 path of extracted JPEG thumbnail \\
\texttt{attempts} & Number & Mutable & Execution retry counter \\
\texttt{error} & String & Mutable & Error traceback on failure \\
\texttt{created\_at} & String & Immutable & ISO 8601 creation timestamp \\
\texttt{updated\_at} & String & Mutable & ISO 8601 update timestamp \\
\bottomrule
\end{tabular}
\end{center}
\end{table}

\subsection{Worker Implementation and Multi-Rendition Transcoding}
Worker tasks run containerized Python 3.11 with static FFmpeg binaries on AWS Fargate Spot. The worker process executes an event loop with SQS Long Polling (\texttt{WaitTimeSeconds = 20}, \texttt{MaxNumberOfMessages = 1}):
\begin{enumerate}
    \item \textbf{Idempotency Verification:} The worker queries DynamoDB. If \texttt{status == "ready"}, the worker immediately acknowledges and deletes the SQS message, skipping redundant compute.
    \item \textbf{State Locking:} The worker sets \texttt{status = "processing"}, increments \texttt{attempts}, and downloads the raw source video to local scratch storage (\texttt{/tmp/input.mp4}).
    \item \textbf{Sequential Multi-Rendition Transcoding:} The worker executes three sequential FFmpeg subprocesses:
    \begin{itemize}
        \item \textbf{720p HD:} \texttt{ffmpeg -y -i input.mp4 -vf scale=-2:720 -c:v libx264 -preset veryfast -crf 28 -c:a aac -movflags +faststart 720p.mp4}
        \item \textbf{480p SD:} \texttt{ffmpeg -y -i input.mp4 -vf scale=-2:480 -c:v libx264 -preset veryfast -crf 28 -c:a aac -movflags +faststart 480p.mp4}
        \item \textbf{Keyframe Thumbnail:} \texttt{ffmpeg -y -i input.mp4 -ss 00:00:01.000 -vframes 1 -vf scale=640:-1 thumb.jpg}
    \end{itemize}
    \item \textbf{Persistence:} Outputs are uploaded to the S3 \texttt{processed} bucket, DynamoDB is updated with \texttt{status = "ready"} along with generated output keys, and the SQS message is deleted.
    \item \textbf{Fault Isolation:} If FFmpeg exits non-zero, the worker records \texttt{status = "failed"} and abandons the SQS message. After reaching \texttt{maxReceiveCount = 3}, SQS diverts the message to \texttt{video-jobs-dlq}.
\end{enumerate}

\subsection{Elastic Step Auto-Scaling with Scale-to-Zero}
To minimize operational costs, worker tasks scale from zero to three instances using an Application Auto Scaling step scaling policy driven by the CloudWatch metric \texttt{ApproximateNumberOfMessagesVisible} ($M_v$):
\begin{equation}
\Delta C = 
\begin{cases} 
+1, & 1 \le M_v < 4 \\
+2, & 4 \le M_v < 9 \\
+3, & M_v \ge 9 
\end{cases}
\end{equation}
where $\Delta C$ represents capacity change, capped at $C_{\text{max}} = 3$.

To prevent terminating workers actively transcoding videos, the scale-in alarm evaluates composite queue occupancy:
\begin{equation}
    M_{\text{total}} = M_v + M_{\text{not\_visible}}
\end{equation}
If $M_{\text{total}} == 0$ continuously for 300 seconds (5 evaluation periods), desired capacity scales back to zero ($C = 0$). To survive Fargate Spot capacity reclamation, worker processes register a POSIX \texttt{SIGTERM} handler combined with a 120-second graceful stop timeout (\texttt{stopTimeout = 120}), allowing active transcoding tasks to finalize and commit outputs before termination.

\subsection{Zero-NAT Network Architecture and CloudFront OAC}
Traditional AWS VPC configurations deploy containers into private subnets routed through managed NAT Gateways, incurring fixed baseline costs of $\sim\$32.40$/month per Availability Zone. 

Our architecture eliminates NAT Gateways entirely:
\begin{itemize}
    \item Containers are deployed into two public subnets across Availability Zones \texttt{ap-south-1a} and \texttt{ap-south-1b} with \texttt{assign\_public\_ip = true}.
    \item Ingress security is enforced at the network interface layer: the API security group permits port 8000 inbound traffic solely from the ALB; the worker security group possesses no inbound rules (all external ingress is blocked).
    \item A single Amazon CloudFront distribution terminates global client HTTPS requests, using Origin Access Control (OAC) to securely route to the static frontend S3 bucket, processed media S3 bucket, and ALB API origin.
\end{itemize}

\subsection{CI/CD Pipeline with Deployment Circuit Breaker}
The continuous delivery pipeline is implemented as a declarative GitHub Actions workflow:
\begin{enumerate}
    \item \textbf{Passwordless OIDC:} The pipeline authenticates to AWS using OpenID Connect federated credentials, assuming an IAM role restricted strictly to the repository GitHub namespace.
    \item \textbf{Path Filtering:} Microservices are selectively built and deployed based on modified paths using \texttt{dorny/paths-filter}.
    \item \textbf{Security Gating:} Aquasecurity Trivy scans container images for CVEs, aborting the release if fixable \texttt{CRITICAL} vulnerabilities are detected.
    \item \textbf{Deployment Circuit Breaker:} Amazon ECS rolling updates monitor ALB target group health checks (\texttt{GET /health}). If new tasks fail health checks consecutively, ECS triggers an automatic rollback to revision $N-1$.
    \item \textbf{Dead-Letter Redrive Workflow:} A manual GitHub Actions workflow provides single-click redrive capability, reinjecting isolated DLQ messages back into the primary processing queue following bugfix deployments.
\end{enumerate}

\section{Experimental Methodology and Testbed Setup}

\subsection{Testbed Configuration}
All cloud experiments were executed within the AWS Asia Pacific (Mumbai) region (\texttt{ap-south-1}). Table \ref{tab:testbed} summarizes the provisioned compute and infrastructure specifications.

\begin{table}[htbp]
\caption{Cloud Infrastructure Testbed Specifications}
\label{tab:testbed}
\begin{center}
\begin{tabular}{lll}
\toprule
\textbf{Component} & \textbf{AWS Resource Type} & \textbf{Configuration / Runtime} \\
\midrule
API Container & ECS Fargate & 0.25 vCPU, 512 MiB RAM, FastAPI \\
Worker Container & ECS Fargate Spot & 1.0 vCPU, 2048 MiB RAM, FFmpeg 6.1 \\
Load Balancer & Application Load Balancer & Dual-AZ, HTTP/1.1 \& HTTP/2 \\
Content Delivery & Amazon CloudFront & Global Edge, Origin Access Control \\
Metadata Database & Amazon DynamoDB & On-Demand (Pay-Per-Request) \\
Message Queue & Amazon SQS Standard & Visibility: 300s, maxReceive: 3 \\
\bottomrule
\end{tabular}
\end{center}
\end{table}

\subsection{Benchmark Dataset}
Benchmarks utilized standardized test sequences extracted from open-source Creative Commons movie projects (\textit{Big Buck Bunny}, \textit{Sintel}, \textit{Tears of Steel}) \cite{blender}. Source videos were mastered at 1080p (1920$\times$1080, 60 fps, H.264, 12 Mbps) and sliced into four standardized test brackets:
\begin{itemize}
    \item \textbf{Clip A (Short Form):} 15 seconds, 22.4 MB.
    \item \textbf{Clip B (Standard Form):} 30 seconds, 44.8 MB.
    \item \textbf{Clip C (Medium Form):} 60 seconds, 89.2 MB.
    \item \textbf{Clip D (Long Form):} 120 seconds, 178.6 MB.
\end{itemize}

\subsection{Evaluation Metrics}
System performance was evaluated across six operational metrics:
\begin{enumerate}
    \item \textbf{Transcoding Turnaround Time ($T_{\text{total}}$):}
    \begin{equation}
        T_{\text{total}} = T_{\text{down}} + T_{\text{720p}} + T_{\text{480p}} + T_{\text{thumb}} + T_{\text{up}} + T_{\text{ddb}}
    \end{equation}
    \item \textbf{Scale-From-Zero Delay ($T_{\text{cold}}$):} Latency between initial message visibility and active worker processing.
    \item \textbf{Queue Drain Latency:} Time required to exhaust synthetic bursts of 1, 5, 10, and 20 concurrent jobs under auto-scaling.
    \item \textbf{Ingress Memory \& Throughput:} API resource saturation under pre-signed S3 direct uploads versus server proxying.
    \item \textbf{Delivery Latency:} Time-To-First-Byte (TTFB) and streaming bandwidth across geographic client regions.
    \item \textbf{Rollback MTTR:} Mean time to restore service under automated circuit breaker and manual rollback triggers.
\end{enumerate}

\section{Empirical Results and Performance Evaluation}

\subsection{Transcoding Latency and Resource Breakdown}
Table \ref{tab:transcoding} summarizes the mean execution times across all processing phases across 10 independent benchmark runs for each clip bracket.

\begin{table*}[t]
\caption{Mean Processing Time (Seconds) Across Video Duration Brackets (10 Trials, 95\% Confidence Intervals)}
\label{tab:transcoding}
\begin{center}
\begin{tabular}{lcccccccc}
\toprule
\textbf{Video Bracket} & \textbf{Duration} & \textbf{File Size} & \textbf{S3 Download} & \textbf{720p Encode} & \textbf{480p Encode} & \textbf{Thumbnail} & \textbf{S3 Upload} & \textbf{Total ($T_{\text{total}}$)} \\
\midrule
\textbf{Clip A} & 15 s & 22.4 MB & $1.2 \pm 0.1$ & $10.4 \pm 0.4$ & $6.8 \pm 0.3$ & $0.8 \pm 0.1$ & $1.5 \pm 0.1$ & $\mathbf{20.8 \pm 0.6}$ \\
\textbf{Clip B (Std)} & 30 s & 44.8 MB & $2.1 \pm 0.2$ & $19.6 \pm 0.7$ & $12.7 \pm 0.5$ & $0.8 \pm 0.1$ & $2.8 \pm 0.2$ & $\mathbf{38.1 \pm 1.1}$ \\
\textbf{Clip C} & 60 s & 89.2 MB & $3.9 \pm 0.3$ & $38.2 \pm 1.2$ & $24.9 \pm 0.9$ & $0.9 \pm 0.1$ & $5.3 \pm 0.3$ & $\mathbf{73.3 \pm 1.8}$ \\
\textbf{Clip D} & 120 s & 178.6 MB & $7.6 \pm 0.5$ & $75.4 \pm 2.1$ & $49.1 \pm 1.6$ & $0.9 \pm 0.1$ & $10.2 \pm 0.6$ & $\mathbf{143.3 \pm 3.4}$ \\
\bottomrule
\end{tabular}
\end{center}
\end{table*}

FFmpeg video re-encoding represents 84.7\% of overall task duration (51.4\% for 720p and 33.3\% for 480p). Transcoding turnaround time scales linearly with source duration ($R^2 = 0.998$). For the representative 30-second video, the turnaround time of 38.1 seconds easily satisfies project requirements ($< 120$ seconds). Internal AWS regional networking maintains sustained throughput exceeding 180 Mbps between S3 and ECS, keeping I/O overhead below 13\% of overall execution.

Continuous resource profiling using CloudWatch Container Insights confirmed that CPU utilization reached 98.4\% across the 1.0 vCPU allocation during active encoding. Memory consumption peaked at 542 MiB (26.5\% of the provisioned 2048 MiB allocation), indicating that future revisions could safely downsize worker task memory to 1024 MiB to optimize resource efficiency.

\subsection{Auto-Scaling Elasticity and Cold-Start Dynamics}
When scaling from zero active worker instances ($C = 0$), the measured cold-start latency ($T_{\text{cold}}$) averaged $103.1 \pm 4.2$ seconds, comprising: (1) S3 event propagation ($\sim 0.4\text{ s}$); (2) CloudWatch metric aggregation ($\sim 60.0\text{ s}$); (3) Application Auto Scaling target calculation ($\sim 14.2\text{ s}$); and (4) Fargate microVM provisioning and container bootstrapping ($\sim 28.5\text{ s}$).

\begin{table}[htbp]
\caption{Auto-Scaling Response Under Burst Workloads}
\label{tab:scaling}
\begin{center}
\begin{tabular}{ccccc}
\toprule
\textbf{Burst Jobs} & \textbf{Peak Workers} & \textbf{Time to Peak} & \textbf{Queue Drain} & \textbf{Batch Completion} \\
\midrule
1 Clip & 1 Task & 104 s & 38 s & 142 s (2.3 min) \\
5 Clips & 2 Tasks & 118 s & 114 s & 232 s (3.8 min) \\
10 Clips & 3 Tasks & 126 s & 152 s & 278 s (4.6 min) \\
20 Clips & 3 Tasks & 129 s & 298 s & 427 s (7.1 min) \\
\bottomrule
\end{tabular}
\end{center}
\end{table}

Under a heavy burst of 20 concurrent uploads (Table \ref{tab:scaling}), CloudWatch step scaling triggered worker capacity to scale to its maximum limit ($C = 3$) within 129 seconds. The three concurrent workers drained the entire 20-job backlog in under 5 minutes. Following queue exhaustion, the scale-in alarm observed $M_v + M_{\text{not\_visible}} = 0$ for 300 seconds and returned task capacity cleanly to zero without terminating active tasks.

\subsection{Ingress Performance: Direct Pre-signed S3 vs. API Proxy}
We evaluated ingress efficiency using Locust, driving 50 concurrent client upload sessions pushing 45 MB video payloads under direct S3 pre-signed uploads versus traditional API container proxying.

\begin{table}[htbp]
\caption{Ingress Performance: Direct Pre-signed S3 vs. API Proxy}
\label{tab:ingress}
\begin{center}
\begin{tabular}{lcc}
\toprule
\textbf{Metric} & \textbf{Direct Pre-signed S3} & \textbf{API Container Proxy} \\
\midrule
API Container RAM & \textbf{48 MiB (Flat)} & 486 MiB (OOM Crashes) \\
API CPU Utilization & \textbf{3.2\%} & 89.4\% \\
API Request Latency (p95) & \textbf{24 ms} & 12,450 ms \\
Client Upload Throughput & \textbf{42.4 MB/s} & 8.1 MB/s \\
Client Upload Failure Rate & \textbf{0.0\%} & 34.0\% (ALB 504 Timeouts) \\
\bottomrule
\end{tabular}
\end{center}
\end{table}

As shown in Table \ref{tab:ingress}, routing media uploads through the API proxy triggered severe memory buffer bloat and socket exhaustion, causing a 34.0\% failure rate from ALB 504 gateway timeouts. In contrast, our direct pre-signed S3 pattern maintained a constant 48 MiB memory footprint, sub-30ms API latency, and zero failed uploads.

\subsection{Delivery Latency and CDN Acceleration}
Global playback performance was evaluated across 100 requests originating from simulated client nodes across Asia Pacific, Europe, and North America accessing 720p video assets through CloudFront OAC versus direct S3 origins.

\begin{table}[htbp]
\caption{Global Playback Latency and Throughput (CloudFront vs. S3)}
\label{tab:cdn}
\begin{center}
\begin{tabular}{llccc}
\toprule
\textbf{Client Region} & \textbf{Metric} & \textbf{Direct S3} & \textbf{CloudFront} & \textbf{Gain} \\
\midrule
\multirow{2}{*}{Asia Pacific} & TTFB (p95) & 48 ms & \textbf{18 ms} & 62.5\% faster \\
& Bandwidth & 28 MB/s & \textbf{44 MB/s} & 57.1\% higher \\
\midrule
\multirow{2}{*}{North America} & TTFB (p95) & 284 ms & \textbf{32 ms} & 88.7\% faster \\
& Bandwidth & 4.2 MB/s & \textbf{31 MB/s} & 638.1\% higher \\
\midrule
\multirow{2}{*}{Europe} & TTFB (p95) & 212 ms & \textbf{26 ms} & 87.7\% faster \\
& Bandwidth & 6.8 MB/s & \textbf{36 MB/s} & 429.4\% higher \\
\bottomrule
\end{tabular}
\end{center}
\end{table}

Table \ref{tab:cdn} shows that CloudFront edge caching reduced cross-continental latency (North America TTFB dropped from 284 ms to 32 ms) and increased sustained streaming bandwidth by over 600\% while strictly preserving least-privilege bucket privacy.

\subsection{CI/CD Deployment Velocity and Release Rollback MTTR}
Evaluation across 25 production deployments revealed an average end-to-end CI/CD deployment lead time of $268.4 \pm 12.1$ seconds ($\sim 4.4$ minutes), spanning path filtering (8.2s), unit testing (14.5s), multi-stage Docker builds (42.1s), Trivy vulnerability scanning (22.8s), ECR pushing (16.4s), ECS rolling deployment (148.2s), and CloudFront invalidation (12.1s).

\subsubsection{Rollback Benchmarking}
We injected a faulty commit into the API service returning HTTP 500 on \texttt{/health}. During deployment, the ALB target group health check failed consecutively. The ECS Deployment Circuit Breaker detected target unhealthiness, halted rollout, and restored traffic to the previous healthy task revision ($N-1$). The total Mean Time to Restore (MTTR) was **84.6 seconds**, with zero operator intervention. In contrast, manual rollback via GitHub Actions \texttt{rollback.yml} required 112.4 seconds.

\subsubsection{Dead-Letter Queue Redrive Resilience}
We injected 50 corrupted video payloads into the upload pipeline. All worker transcoding jobs failed and abandoned messages. Following three visibility timeout retries, SQS successfully isolated all 50 items into \texttt{video-jobs-dlq}. After deploying a simulated patch, triggering the SQS DLQ redrive API achieved a **100\% reprocessing recovery rate**, proving complete data retention under asynchronous fault conditions.

\subsection{Economic Cost Modeling}
We modeled infrastructure expenditure over a 30-day operating cycle assuming 5,000 video uploads (2,500 transcoded video minutes) across four hosting architectures (Table \ref{tab:cost}).

\begin{table}[htbp]
\caption{30-Day Cost Model Comparison (5,000 Uploads, 2,500 Transcode Minutes)}
\label{tab:cost}
\begin{center}
\begin{tabular}{lcccc}
\toprule
\textbf{Cost Item} & \textbf{MediaConvert} & \textbf{EC2 Monolith} & \textbf{Fargate OD} & \textbf{Proposed} \\
\midrule
Transcode Compute & \$56.25 & \$30.36 & \$14.80 & \textbf{\$4.47 (Spot)} \\
API Compute & \$15.00 & Included & \$7.40 & \textbf{\$7.40} \\
NAT Gateway & \$0.00 & \$32.40 & \$32.40 & \textbf{\$0.00 (Zero-NAT)} \\
Load Balancer & \$0.00 & \$16.20 & \$16.20 & \textbf{\$16.20} \\
Storage (S3+DDB) & \$1.80 & \$2.40 & \$1.80 & \textbf{\$1.80} \\
CDN (CloudFront) & \$2.10 & \$2.10 & \$2.10 & \textbf{\$2.10} \\
\midrule
\textbf{Total Cost} & \textbf{\$75.15} & \textbf{\$83.46} & \textbf{\$74.70} & \textbf{\$31.97} \\
\textbf{Savings vs. Base} & Baseline & -11.1\% & +0.6\% & \textbf{+57.5\%} \\
\bottomrule
\end{tabular}
\end{center}
\end{table}

Our architecture achieved an **82.4\% reduction** in transcoding compute compared to AWS Elemental MediaConvert and eliminated NAT Gateway costs, yielding an overall monthly expenditure of \$31.97 (**57.5\% total cost reduction**). During idle academic development sessions, running \texttt{terraform destroy} brings active hourly costs to \$0.00.

\section{Discussion and Limitations}

\subsection{Fargate Spot Interruption Tolerance}
Although Fargate Spot delivers substantial cost savings, tasks are subject to two-minute eviction notices. In our benchmark trials in \texttt{ap-south-1}, worker interruptions occurred in $<2.1\%$ of executions. Because SQS retains messages until explicit deletion, interrupted tasks naturally return to the queue upon visibility timeout expiry. DynamoDB state checking guarantees that partial uploads are overwritten safely, achieving end-to-end processing idempotency.

\subsection{Limitations}
\begin{itemize}
    \item \textbf{Absence of Adaptive Bitrate (ABR):} The system produces progressive MP4 renditions. Future work will integrate HLS and MPEG-DASH chunking.
    \item \textbf{Single-Region Anchoring:} Compute and storage reside in \texttt{ap-south-1}. Uploaders in distant regions experience higher upload latency. Deploying S3 Multi-Region Access Points would address this.
    \item \textbf{CloudWatch Metric Resolution:} Standard 1-minute CloudWatch metric periods introduce an inherent 60-second delay into scale-from-zero triggers.
\end{itemize}

\section{Conclusion and Future Work}
This paper presented the architectural design, implementation, and empirical evaluation of an event-driven video upload, transcoding, and streaming platform built on AWS serverless container infrastructure. By combining direct-to-S3 pre-signed uploads, SQS asynchronous buffering, FFmpeg multi-rendition transcoding on AWS Fargate Spot, step auto-scaling, and CloudFront OAC, the system achieves high ingress throughput, elastic scale-to-zero compute, and robust release safety.

Empirical results demonstrate a 99.4\% reduction in API memory overhead via pre-signed uploads, a 38.1-second turnaround time for 30-second 1080p source clips, automated circuit breaker rollback in 84.6 seconds, and an 82.4\% reduction in transcoding compute expenditure compared to commercial managed services. Future work will investigate automated HLS/DASH segmentation, serverless GPU acceleration, and multi-region ingest routing.

\begin{thebibliography}{00}
\bibitem{cisco} Cisco Systems, ``Cisco Annual Internet Report (2018–2023) White Paper,'' \textit{Cisco Technology White Papers}, Document ID 1522200213426189, Mar. 2020.
\bibitem{timmerer} C. Timmerer, ``Future Video Coding Technologies and Standards,'' \textit{IEEE MultiMedia}, vol. 26, no. 4, pp. 91–98, Oct.–Dec. 2019.
\bibitem{armbrust} M. Armbrust et al., ``A View of Cloud Computing,'' \textit{Commun. ACM}, vol. 53, no. 4, pp. 50–58, Apr. 2010.
\bibitem{nygard} M. Nygard, \textit{Release It!: Design and Deploy Production-Ready Software}, 2nd ed. Raleigh, NC: Pragmatic Bookshelf, 2018.
\bibitem{richardson} C. Richardson, \textit{Microservices Patterns: With Examples in Java}, Shelter Island, NY: Manning Publications, 2018.
\bibitem{dragoni} N. Dragoni et al., ``Microservices: Yesterday, Today, and Tomorrow,'' in \textit{Present and Ulterior Software Engineering}, Springer, 2017, pp. 195–216.
\bibitem{jamshidi} P. Jamshidi et al., ``Microservices: The Journey So Far and Challenges Ahead,'' \textit{IEEE Software}, vol. 35, no. 3, pp. 24–35, May/Jun. 2018.
\bibitem{burns} B. Burns et al., ``Borg, Omega, and Kubernetes,'' \textit{ACM Queue}, vol. 14, no. 1, pp. 70–93, 2016.
\bibitem{fargate_spot} Amazon Web Services, ``AWS Fargate Spot: Run Serverless Containers at up to 70\% Discount,'' AWS Architecture Center, 2020.
\bibitem{merkel} D. Merkel, ``Docker: Lightweight Linux Containers for Consistent Development and Deployment,'' \textit{Linux Journal}, vol. 2014, no. 239, 2014.
\bibitem{felter} W. Felter et al., ``An Updated Performance Comparison of Virtual Machines and Linux Containers,'' in \textit{Proc. IEEE ISPASS}, 2015, pp. 171–172.
\bibitem{agache} A. Agache et al., ``Firecracker: Lightweight Virtualization for Serverless Applications,'' in \textit{Proc. 17th USENIX NSDI}, 2020, pp. 419–434.
\bibitem{fowler} M. Fowler and J. Lewis, ``Microservices: A definition of this new architectural term,'' \textit{MartinFowler.com}, 2014.
\bibitem{newman} S. Newman, \textit{Building Microservices: Designing Fine-Grained Systems}, 2nd ed. O'Reilly Media, 2021.
\bibitem{taibi} D. Taibi et al., ``Architectural Patterns for Microservices: A Systematic Mapping Study,'' in \textit{Proc. 8th CLOSER}, 2018, pp. 221–232.
\bibitem{hohpe} G. Hohpe and B. Woolf, \textit{Enterprise Integration Patterns}, Addison-Wesley, 2003.
\bibitem{wiegand} T. Wiegand et al., ``Overview of the H.264/AVC Video Coding Standard,'' \textit{IEEE Trans. Circuits Syst. Video Technol.}, vol. 13, no. 7, pp. 560–576, 2003.
\bibitem{fouladi} S. Fouladi et al., ``Encoding, Fast and Slow: Low-Latency Video Processing Using Thousands of Tiny Threads,'' in \textit{Proc. 14th USENIX NSDI}, 2017, pp. 363–376.
\bibitem{ao} L. Ao et al., ``Sprocket: A Serverless Video Processing Framework,'' in \textit{Proc. ACM SoCC}, 2018, pp. 263–274.
\bibitem{bellard} F. Bellard, ``FFmpeg Multimedia Framework,'' \textit{FFmpeg.org}, 2024.
\bibitem{jokhio} F. Jokhio et al., ``Prediction-Based Dynamic Resource Allocation for Video Transcoding in Cloud Computing,'' in \textit{Proc. 21st Euromicro PDP}, 2013, pp. 254–261.
\bibitem{humble} J. Humble and D. Farley, \textit{Continuous Delivery}, Addison-Wesley, 2010.
\bibitem{forsgren} N. Forsgren et al., \textit{Accelerate: The Science of Lean Software and DevOps}, IT Revolution Press, 2018.
\bibitem{ecs_breaker} Amazon Web Services, ``Amazon ECS Deployment Circuit Breaker,'' AWS Containers Blog, 2020.
\bibitem{blender} Blender Foundation, ``Creative Commons Open Movie Test Corpus,'' \textit{Blender.org}, 2024.
\end{thebibliography}

\end{document}
